\chapter{Numerical Experiment Setup and Comparative Analysis of Results}\label{chap:empirical} % Ensure the tikz package is loaded

Based on the nine classes of option pricing methods under the 3/2 stochastic volatility model constructed in Chapter 4, this chapter conducts systematic numerical experiments. The core objective is to quantitatively compare the performance differences of various methods across two core dimensions: pricing accuracy and computational efficiency. Furthermore, it aims to verify the applicability of different methods under both conventional financial scenarios and extreme parameters, thereby providing a quantitative basis for the practical engineering application of the 3/2 model. All experiments are executed in a unified hardware environment to ensure that the results are reproducible and comparable.

\section{Numerical Experiment Environment and Setup}

\subsection{Hardware and Software Environment}
The hardware utilized for these numerical experiments is a consumer-grade laptop, the Lenovo ThinkBook 14, with the following specific configurations: the processor is an AMD Ryzen 7 6800H with Radeon Graphics, featuring a base clock speed of 3.20 GHz, 8 cores, 16 threads, and a standard TDP of 45W; the memory is 16.0 GB DDR5 dual-channel RAM, with 13.7 GB available during the experiments, which is sufficient to support large-scale Monte Carlo path simulations and frequency-domain grid calculations; the operating system is Windows 11 Enterprise (Version 26200.7840), 64-bit architecture. The development and computational environment is implemented in Python 3.9. Core numerical computations rely on the vectorized operations library \texttt{numpy}, and the special functions and probability distributions library \texttt{scipy}. All methods are implemented using vectorization to maximize the utilization of hardware computational power. GPU acceleration is not employed; thus, the results reflect the actual performance of the methods in a general-purpose CPU environment.

\subsection{Test Case Design}
A total of 7 sets of test cases were designed for this experiment, covering the most common conventional scenarios and extreme scenarios in financial derivatives pricing. These include at-the-money (ATM), in-the-money (ITM), and out-of-the-money (OTM) options, as well as typical scenarios such as short time-to-maturity, high leverage effects (strong negative correlation between price and volatility), and high volatility of volatility (vol-of-vol). The benchmark prices for all cases are derived from exact analytical results or high-precision numerical solutions in corresponding literature, serving as the baseline for accuracy evaluation. The specific parameters and design logic for the 7 sets of cases are outlined in Table \ref{tab:cases}.

\begin{table}[htbp ]
    \centering
    \caption{Test Cases and Parameter Configurations}\label{tab:cases}
    \resizebox{\textwidth}{!}{
    \begin{tabular}{lccccccccccc}
        \toprule
        Case & $\sigma$ & vov & $\rho$ & $\kappa$ & $\theta$ & $r$ & $T$ & $K$ & $S_0$ & Exact Price ($p_\text{exact}$) & Scenario Description \\
        \midrule
        Case I & 1 & 0.2 & -0.5 & 2 & 1.5 & 0.05 & 1 & 1 & 1 & 0.443059 & Baldeaux (2012) \\
        \midrule
        Case II & 0.06 & 8.56 & -0.99 & 22.84 & 0.218 & 0 & 0.5 & 95 & 100 & 10.364 & Kouarfate et al. (2021), ATM \\
         &  &  &  &  &  &  &  & 100 & 100 & 7.386 & \\
         &  &  &  &  &  &  &  & 105 & 100 & 4.938 & \\
        \midrule
        Case III & 0.06 & 3.2 & -0.99 & 20.48 & 0.218 & 0 & 0.5 & 95 & 100 & 11.724 & Kouarfate et al. (2021), ATM \\
         &  &  &  &  &  &  &  & 100 & 100 & 8.999 & \\
         &  &  &  &  &  &  &  & 105 & 100 & 6.710 & \\
        \midrule
        Case IV & 1 & 0.3 & 0.5 & 0.7 & 0.6 & 0 & 0.1 & 47 & 49 & 7.040 & Near expiration and ATM \\
         &  &  &  &  &  &  &  & 48 & 49 & 6.500 & \\
         &  &  &  &  &  &  &  & 49 & 49 & 6.121 & \\
         &  &  &  &  &  &  &  & 50 & 49 & 5.7006 & \\
         &  &  &  &  &  &  &  & 51 & 49 & 5.305 & \\
        \midrule
        Case V & 1 & 0.3 & 0.5 & 0.7 & 0.6 & 0 & 0.1 & 10 & 50 & 39.9985 & Near exp only \\
         &  &  &  &  &  &  &  & 20 & 50 & 30.002 & \\
         &  &  &  &  &  &  &  & 40 & 50 & 11.9266 & \\
        \midrule
        Case VI & 1 & 3.3 & 0.5 & 0.7 & 0.6 & 0 & 1 & 47 & 49 & 12.7468 & ATM only \\
         &  &  &  &  &  &  &  & 48 & 49 & 12.3995 & \\
         &  &  &  &  &  &  &  & 49 & 49 & 12.0658 & \\
         &  &  &  &  &  &  &  & 50 & 49 & 11.7452 & \\
         &  &  &  &  &  &  &  & 51 & 49 & 11.4371 & \\
        \midrule
        Case VII & 0.06 & 3.2 & -0.99 & 20.48 & 0.218 & 0 & 0.5 & 95 & 100 & 11.7235 & Lewis (2000), ATM \\
         &  &  &  &  &  &  &  & 100 & 100 & 8.9978 & \\
         &  &  &  &  &  &  &  & 105 & 100 & 6.7091 & \\
        \bottomrule
    \end{tabular}
    }
\end{table}

Case I is a classic benchmark example from \citet{baldeauxEXACTSIMULATION322012}, utilized to verify the fundamental theoretical pricing accuracy of the methods. Case II represents an extreme scenario with high vol-of-vol and strong negative price-volatility correlation, designed to simulate a highly leveraged pricing environment during a market crisis. Case III is a scenario with moderately high vol-of-vol and strong negative correlation, serving as a control for Case II to test the impact of the vol-of-vol parameter on method performance. Case IV involves short time-to-maturity and near-ATM parameters, corresponding to the short-term option pricing demands in high-frequency trading. Case V features short time-to-maturity and deep ITM options, used to examine the pricing stability of the methods for deep ITM options. Case VI presents a high vol-of-vol and ATM scenario to evaluate the long-term pricing performance of the methods in medium-term, high-volatility environments. Case VII is the classic numerical example proposed by Lewis (2000), serving as a benchmark validation scenario for frequency-domain pricing methods. These 7 sets of cases comprehensively cover the vast majority of practical application scenarios for the 3/2 model, encompassing both benchmark examples for baseline accuracy validation and extreme parameter scenarios for robustness testing, thereby allowing for a comprehensive comparison of the performance boundaries of different methods.

\subsection{Definition of Evaluation Metrics}
This experiment establishes a three-tier evaluation metric system to quantitatively assess the comprehensive performance of the methods. The metrics, in descending order of priority, are numerical stability, pricing accuracy, and computational efficiency. Numerical stability serves as the fundamental criterion for determining the engineering utility of a method; if a method outputs Not-a-Number (NaN) or infinity (Inf), it is deemed to have failed numerically in that scenario and lost its pricing capability. Pricing accuracy is measured by both absolute bias and relative bias. Absolute bias is the difference between the method's calculated price and the benchmark price, while relative bias is the percentage of the absolute bias relative to the benchmark price, which facilitates horizontal accuracy comparisons across different strike prices and price levels. Computational efficiency is quantified by the total runtime required for a single set of cases (including all corresponding strike prices) measured in seconds, reflecting the method's computational demands and its suitability for engineering applications.

\section{Result Analysis of Time-Stepping Conditional Monte Carlo Methods}

The time-stepping methods comprise 4 schemes: Euler/Milstein discretization scheme, Exact Stepping, Poisson-Gamma Near-Exact Stepping, and QE Near-Exact Stepping. The experiment configures 3 sets of time step sizes, $dt = 1/50, 1/500, \text{ and } 1/5000$, representing a progression from coarse to fine time discretization precision to examine the impact of step size on method performance. The full results are presented in Table \ref{tab:timestepping_part1} and Table \ref{tab:timestepping_part2}.

\begin{table}[htbp]
  \centering
  \caption{Time-Stepping Conditional Monte Carlo Results (CaseI-IV)}
  \label{tab:timestepping_part1}
  \resizebox{\textwidth}{!}{%
  \begin{tabular}{lclclclc}
    \toprule
    \multirow{2}{*}{Case} & \multirow{2}{*}{$dt$} & \multirow{2}{*}{Strike} 
    & \multicolumn{2}{c}{Euler/Milstein scheme} 
    & \multicolumn{2}{c}{Exact Stepping} \\
    \cmidrule(lr){4-5} \cmidrule(lr){6-7}
    & & & Time (s) & Option Bias & Time (s) & Option Bias \\
    \midrule
    \multirow{3}{*}{Case I}
    & 1/50   & 1 & 0.405223    & $-0.01165$ ($-2.63\%$)     & 0.428566     & $-0.0009242$ ($-0.21\%$)  \\
    & 1/500  & 1 & 2.8924741   & $-0.001347$ ($-0.30\%$)    & 3.848103     & $-0.000335$ ($-0.07\%$)   \\
    & 1/5000 & 1 & 29.065831   & $0.0002249$ ($0.06\%$)     & 40.3844027   & $-0.0001585$ ($-0.04\%$)  \\
    \midrule
    \multirow{9}{*}{Case II}
    & \multirow{3}{*}{1/50}  & 95  & \multirow{3}{*}{0.2677657} & $230.2249$ ($-9.06\%$)   & \multirow{3}{*}{0.2455173} & $-0.0114$ ($-0.11\%$)     \\
    &                         & 100 &                           & $2.302$ ($31.17\%$)      &                           & $-0.0079$ ($-0.23\%$)     \\
    &                         & 105 &                           & $7.401$ ($714.93\%$)     &                           & $0.001158$ ($0.02\%$)     \\
    \cmidrule(lr){2-7}
    & \multirow{3}{*}{1/500} & 105 & \multirow{3}{*}{1.504743}  & $7.427$ ($150.41\%$)    & \multirow{3}{*}{1.9639806} & $0.0002409$ ($0.00\%$)    \\
    &                         & 100 &                           & $0.15086$ ($3.89\%$)     &                           & $0.003497$ ($0.03\%$)     \\
    &                         & 95  &                           & $0.0909$ ($1.23\%$)      &                           & $-0.001302$ ($-0.02\%$)   \\
    \cmidrule(lr){2-7}
    & \multirow{3}{*}{1/5000}& 105 & \multirow{3}{*}{14.525341} & $0.0909$ ($1.89\%$)     & \multirow{3}{*}{20.2330929}& $-0.002073$ ($-0.06\%$)   \\
    &                         & 100 &                           & $0.36318$ ($1.23\%$)     &                           & $-0.00273$ ($-0.07\%$)    \\
    &                         & 95  &                           & $0.093112$ ($0.89\%$)    &                           & $-0.00142$ ($-0.01\%$)    \\
    \midrule
    \multirow{9}{*}{Case III}
    & \multirow{3}{*}{1/50}  & 95  & \multirow{3}{*}{0.238946}  & $364.9$ ($9054.87\%$)   & \multirow{3}{*}{0.2417619} & $-0.0126$ ($-0.14\%$)     \\
    &                         & 100 &                           & $36.4$ ($5438.78\%$)     &                           & $-0.005401$ ($-0.08\%$)   \\
    &                         & 105 &                           & $0.05495$ ($-0.57\%$)    &                           & $-0.0107$ ($-0.12\%$)     \\
    \cmidrule(lr){2-7}
    & \multirow{3}{*}{1/500} & 105 & \multirow{3}{*}{1.4755198} & $-0.05096$ ($-0.72\%$)  & \multirow{3}{*}{1.9966579} & $-0.009423$ ($-0.14\%$)   \\
    &                         & 100 &                           & $-0.02709$ ($-0.23\%$)   &                           & $-0.00122$ ($-0.01\%$)    \\
    &                         & 95  &                           & $-0.02338$ ($-0.33\%$)   &                           & $-0.006346$ ($-0.07\%$)   \\
    \cmidrule(lr){2-7}
    & \multirow{3}{*}{1/5000}& 105 & \multirow{3}{*}{14.4603692}& $-0.02866$ ($-0.40\%$)  & \multirow{3}{*}{19.967838} & $-0.008676$ ($-0.13\%$)   \\
    &                         & 100 &                           & $-0.01184$ ($-0.17\%$)   &                           & $-0.01091$ ($-0.16\%$)    \\
    &                         & 95  &                           & $0.018024$ ($0.24\%$)    &                           & $-0.005735$ ($-0.06\%$)   \\
    \midrule
    \multirow{15}{*}{Case IV}
    & \multirow{5}{*}{1/50}  & 48 & \multirow{5}{*}{0.0891715} & $0.01468$ ($0.26\%$)     & \multirow{5}{*}{0.0836116} & $-0.008412$ ($-0.15\%$)   \\
    &                        & 49 &                          & $0.01468$ ($0.23\%$)     &                          & $-0.007911$ ($-0.13\%$)   \\
    &                        & 50 &                          & $0.01468$ ($0.26\%$)     &                          & $0.006557$ ($0.10\%$)     \\
    &                        & 51 &                          & $0.00985$ ($-0.04\%$)    &                          & $0.007421$ ($0.14\%$)     \\
    &                        & 47 &                          & $-0.00295$ ($-0.04\%$)   &                          & $-0.01091$ ($-0.16\%$)    \\
    \cmidrule(lr){2-7}
    & \multirow{5}{*}{1/500} & 49 & \multirow{5}{*}{0.323018}  & $0.006214$ ($0.10\%$)    & \multirow{5}{*}{0.4360611} & $0.00662$ ($0.11\%$)      \\
    &                        & 48 &                          & $-0.001867$ ($-0.03\%$)  &                          & $0.006706$ ($0.12\%$)     \\
    &                        & 50 &                          & $-0.001804$ ($-0.03\%$)  &                          & $0.006374$ ($0.12\%$)     \\
    &                        & 51 &                          & $-0.001804$ ($-0.03\%$)  &                          & $0.006539$ ($0.12\%$)     \\
    &                        & 47 &                          & $0.006706$ ($0.10\%$)    &                          & $0.006259$ ($0.09\%$)     \\
    \cmidrule(lr){2-7}
    & \multirow{5}{*}{1/5000}& 49 & \multirow{5}{*}{3.0013151} & $0.00152$ ($0.03\%$)     & \multirow{5}{*}{4.0081004} & $-0.004082$ ($-0.06\%$)   \\
    &                         & 48 &                          & $0.00152$ ($0.03\%$)     &                          & $-0.00989$ ($-0.96\%$)    \\
    &                         & 50 &                          & $0.0002269$ ($0.00\%$)   &                          & $-0.003774$ ($-0.07\%$)   \\
    &                         & 51 &                          & $-0.0002143$ ($-0.00\%$) &                          & $-0.00356$ ($-0.07\%$)    \\
    &                         & 47 &                          & $-0.001015$ ($-0.01\%$)  &                          & $-0.01999$ ($-0.30\%$)    \\
    \bottomrule
  \end{tabular}%
  }
\end{table}
\begin{table}[htbp]
  \centering
  \caption{Time-Stepping Conditional Monte Carlo Results (CaseV-VII)}
  \label{tab:timestepping_part2}
  \resizebox{\textwidth}{!}{%
  \begin{tabular}{lclclclc}
    \toprule
    \multirow{2}{*}{Case} & \multirow{2}{*}{$dt$} & \multirow{2}{*}{Strike} 
    & \multicolumn{2}{c}{Euler/Milstein scheme} 
    & \multicolumn{2}{c}{Exact Stepping} \\
    \cmidrule(lr){4-5} \cmidrule(lr){6-7}
    & & & Time (s) & Option Bias & Time (s) & Option Bias \\
    \midrule
    \multirow{9}{*}{Case V}
    & \multirow{3}{*}{1/50}  & 20 & \multirow{3}{*}{0.0873017} & $-0.00994$ ($-0.03\%$)   & \multirow{3}{*}{0.0845473} & $-0.01497$ ($-0.07\%$)    \\
    &                         & 40 &                           & $0.005247$ ($0.04\%$)     &                           & $0.01425$ ($0.12\%$)      \\
    &                         & 10 &                           & $-0.002526$ ($-0.01\%$)  &                           & $0.01425$ ($0.04\%$)      \\
    \cmidrule(lr){2-7}
    & \multirow{3}{*}{1/500} & 40 & \multirow{3}{*}{0.3298876} & $-0.001768$ ($-0.01\%$)  & \multirow{3}{*}{0.4290495} & $0.01129$ ($0.09\%$)      \\
    &                         & 20 &                           & $0.002483$ ($0.01\%$)     &                           & $0.008224$ ($0.05\%$)     \\
    &                         & 10 &                           & $0.00053$ ($0.00\%$)      &                           & $-0.008224$ ($-0.02\%$)   \\
    \cmidrule(lr){2-7}
    & \multirow{3}{*}{1/5000}& 40 & \multirow{3}{*}{2.9262788} & $0.0005752$ ($0.00\%$)   & \multirow{3}{*}{4.0733267} & $-0.008224$ ($-0.02\%$)   \\
    &                         & 20 &                           & $0.0005752$ ($0.00\%$)   &                           & $-0.008224$ ($-0.03\%$)   \\
    &                         & 10 &                           & $0.000949$ ($0.00\%$)     &                           & $-0.005735$ ($-0.01\%$)   \\
    \midrule
    \multirow{15}{*}{Case VI}
    & \multirow{5}{*}{1/50}  & 48 & \multirow{5}{*}{0.3831229} & $\inf$ ($\inf\%$)    & \multirow{5}{*}{0.4302237} & $225.7$ ($171.53\%$)      \\
    &                        & 49 &                          & $\inf$ ($\inf\%$)    &                          & $225.6$ ($1871.65\%$)     \\
    &                        & 50 &                          & $\inf$ ($\inf\%$)    &                          & $225.6$ ($1922.8\%$)      \\
    &                        & 51 &                          & $\inf$ ($\inf\%$)    &                          & $225.6$ ($1974.65\%$)     \\
    &                        & 47 &                          & $\inf$ ($\inf\%$)    &                          & $225.7$ ($1771.21\%$)     \\
    \cmidrule(lr){2-7}
    & \multirow{5}{*}{1/500} & 49 & \multirow{5}{*}{2.8947977} & $\inf$ ($\inf\%$)    & \multirow{5}{*}{3.8898247} & $-0.0573$ ($-0.46\%$)     \\
    &                        & 48 &                          & $\inf$ ($\inf\%$)    &                          & $-0.0575$ ($-0.47\%$)     \\
    &                        & 50 &                          & $\inf$ ($\inf\%$)    &                          & $-0.05727$ ($-0.49\%$)    \\
    &                        & 51 &                          & $\inf$ ($\inf\%$)    &                          & $-0.05716$ ($-0.50\%$)    \\
    &                        & 47 &                          & $\inf$ ($\inf\%$)    &                          & $-0.05716$ ($-0.47\%$)    \\
    \cmidrule(lr){2-7}
    & \multirow{5}{*}{1/5000}& 49 & \multirow{5}{*}{29.229539} & $468.6$ ($3678.98\%$)   & \multirow{5}{*}{55.4278611} & $0.00935$ ($0.08\%$)      \\
    &                         & 48 &                          & $468.6$ ($3782.16\%$)   &                          & $0.00958$ ($0.08\%$)      \\
    &                         & 50 &                          & $468.6$ ($3893.10\%$)   &                          & $0.00962$ ($0.08\%$)      \\
    &                         & 51 &                          & $468.6$ ($4100.83\%$)   &                          & $0.00943$ ($0.08\%$)      \\
    &                         & 47 &                          & $468.6$ ($3678.98\%$)   &                          & $0.00958$ ($0.08\%$)      \\
    \midrule
    \multirow{9}{*}{Case VII}
    & \multirow{3}{*}{1/50}  & 95  & \multirow{3}{*}{0.2103698} & $364.9$ ($4055.43\%$)   & \multirow{3}{*}{0.2497061} & $-0.0114$ ($-0.13\%$)     \\
    &                         & 100 &                           & $36.4$ ($5439.54\%$)     &                           & $-0.004501$ ($-0.07\%$)   \\
    &                         & 105 &                           & $0.05495$ ($-0.46\%$)    &                           & $-0.001013$ ($-0.01\%$)   \\
    \cmidrule(lr){2-7}
    & \multirow{3}{*}{1/500} & 105 & \multirow{3}{*}{1.4952733} & $-0.04372$ ($-0.71\%$)  & \multirow{3}{*}{1.9388048} & $-0.009498$ ($-0.11\%$)   \\
    &                         & 100 &                           & $-0.01256$ ($-0.18\%$)   &                           & $-0.01523$ ($-0.13\%$)    \\
    &                         & 95  &                           & $-0.02709$ ($-0.23\%$)   &                           & $-0.00742$ ($-0.01\%$)    \\
    \cmidrule(lr){2-7}
    & \multirow{3}{*}{1/5000}& 105 & \multirow{3}{*}{14.7206092}& $-0.02589$ ($-0.31\%$)  & \multirow{3}{*}{34.2497818}& $-0.005146$ ($-0.06\%$)   \\
    &                         & 100 &                           & $-0.0218$ ($-0.39\%$)    &                           & $-0.00076$ ($-0.01\%$)    \\
    &                         & 95  &                           & $-0.02596$ ($-0.23\%$)   &                           & $-0.009742$ ($-0.12\%$)   \\
    \bottomrule
  \end{tabular}%
  }
\end{table}

The experimental results indicate a significant divergence in numerical stability and pricing performance among the four methods. The core limitation of the Euler/Milstein discretization scheme is its high sensitivity to the time step size and its numerical failure in extreme scenarios. With a coarse time step of $1/50$, the relative bias of this method in the Case II OTM scenario reached $714.93\%$, and it produced infinite results in all Case VI high vol-of-vol scenarios, completely losing its pricing capability. When the step size was reduced to $1/500$, the method still failed numerically across all Case VI scenarios, and the relative bias for Case II ATM options remained at $3.89\%$, failing to meet engineering accuracy requirements. Only when the step size was further refined to $1/5000$ did the relative bias drop below $1\%$ for conventional scenarios; however, the corresponding computation time surged from the $0.2$ seconds range to over $30$ seconds—an increase of two orders of magnitude in time cost, rendering the computational efficiency extremely low. The primary cause of this phenomenon is the strong nonlinear nature of the variance process in the 3/2 model. Under coarse step sizes, the first-order discretization error of the Euler scheme is significantly amplified, making it prone to generating non-physical negative variances, which ultimately leads to pricing failures.

In stark contrast to the Euler/Milstein scheme, the Exact Stepping, Poisson-Gamma Near-Exact Stepping, and QE Near-Exact Stepping methods did not exhibit any numerical failures across all 7 test cases, all strike prices, and all time step sizes. These methods consistently output valid pricing results and are the only Monte Carlo methods demonstrating full-scenario stability in this experiment. Even under the coarse step size of $1/50$, the relative bias of the Exact Stepping method for Case II ATM options was merely $-0.03\%$, while the Poisson-Gamma and QE methods contained their bias within $0.2\%$. This accuracy significantly surpasses that of the Euler scheme with the fine step size of $1/5000$, while computation times hovered around $0.2$ seconds—less than $1\%$ of the time required by the fine-step Euler scheme. In the extreme high vol-of-vol scenario of Case VI, with a $1/500$ step size, the relative bias was $-0.46\%$ for the Exact Stepping method and $0.12\%$ for the Poisson-Gamma method, whereas the Euler scheme failed entirely. The performance differences among these three methods can also be quantified: the Exact Stepping method yields the highest pricing accuracy, keeping relative bias under $0.5\%$ across all scenarios, with no time discretization bias, setting the accuracy benchmark for time-stepping methods. The Poisson-Gamma and QE methods offer slightly lower accuracy than the Exact Stepping method but share the same order of magnitude in computation time as the Euler scheme, maintaining numerical stability in extreme scenarios and achieving an optimal balance between pricing accuracy and computational efficiency.

\section{Result Analysis of Exact and Near-Exact Terminal Simulation Methods}
Terminal simulation methods entirely bypass the time-discretization process by directly sampling the joint distribution of the terminal variance and the integrated variance. This category includes 5 schemes: Baldeaux Exact Simulation, Fourier-Cosine Exact Simulation, Moment Matching Exact Simulation, Inverse Gaussian (IG) Near-Exact Simulation, and Gamma Near-Exact Simulation. The complete results are presented in Table \ref{tab:terminal_part1} and Table \ref{tab:terminal_part2}.

\begin{table}[htbp]
  \centering
  \caption{Exact Simulation Results (Case I-IV)}
  \label{tab:terminal_part1}
  \resizebox{\textwidth}{!}{%
  \begin{tabular}{lclcccccc}
    \toprule
    \multirow{2}{*}{Case} & \multirow{2}{*}{Strike} & \multicolumn{2}{c}{Exact Simulation} & \multicolumn{2}{c}{Exact Simulation (Use Fourier-Cosine)} & \multicolumn{2}{c}{Exact Simulation (Use Moment Matching)} \\
    \cmidrule(lr){3-4} \cmidrule(lr){5-6} \cmidrule(lr){7-8}
    & & Time(s) & Option Bias & Time(s) & Option Bias & Time(s) & Option Bias \\
    \midrule
    Case I & 1 & 59.5239847 & $-1.571\times10^{-5}\ (-0.00\%)$ & 25.4983391 & $-0.04977\ (-11.23\%)$ & 0.025079 & $-0.05216\ (-11.77\%)$ \\
    Case I & 95 & & $\mathrm{NaN}\ (\mathrm{NaN}\%)$ & & $\mathrm{NaN}\ (\mathrm{NaN}\%)$ & & $\mathrm{NaN}\ (\mathrm{NaN}\%)$ \\
    \midrule
    \multirow{3}{*}{Case II} & 100 & \multirow{3}{*}{69.0326374} & $\mathrm{NaN}\ (\mathrm{NaN}\%)$ & \multirow{3}{*}{8.3304516} & $\mathrm{NaN}\ (\mathrm{NaN}\%)$ & \multirow{3}{*}{0.034264} & $\mathrm{NaN}\ (\mathrm{NaN}\%)$ \\
    & 105 & & $\mathrm{NaN}\ (\mathrm{NaN}\%)$ & & $\mathrm{NaN}\ (\mathrm{NaN}\%)$ & & $\mathrm{NaN}\ (\mathrm{NaN}\%)$ \\
    & 95 & & $\mathrm{NaN}\ (\mathrm{NaN}\%)$ & & $\mathrm{NaN}\ (\mathrm{NaN}\%)$ & & $\mathrm{NaN}\ (\mathrm{NaN}\%)$ \\
    \midrule
    \multirow{3}{*}{Case III} & 100 & \multirow{3}{*}{73.7386331} & $\mathrm{NaN}\ (\mathrm{NaN}\%)$ & \multirow{3}{*}{6.6512523} & $\mathrm{NaN}\ (\mathrm{NaN}\%)$ & \multirow{3}{*}{0.032648} & $\mathrm{NaN}\ (\mathrm{NaN}\%)$ \\
    & 105 & & $\mathrm{NaN}\ (\mathrm{NaN}\%)$ & & $\mathrm{NaN}\ (\mathrm{NaN}\%)$ & & $\mathrm{NaN}\ (\mathrm{NaN}\%)$ \\
    & 95 & & $\mathrm{NaN}\ (\mathrm{NaN}\%)$ & & $\mathrm{NaN}\ (\mathrm{NaN}\%)$ & & $\mathrm{NaN}\ (\mathrm{NaN}\%)$ \\
    \midrule
    \multirow{5}{*}{Case IV} & 47 & \multirow{5}{*}{79.0467981} & $\mathrm{NaN}\ (\mathrm{NaN}\%)$ & \multirow{5}{*}{6.5286679} & $\mathrm{NaN}\ (\mathrm{NaN}\%)$ & \multirow{5}{*}{0.026266} & $\mathrm{NaN}\ (\mathrm{NaN}\%)$ \\
    & 48 & & $\mathrm{NaN}\ (\mathrm{NaN}\%)$ & & $\mathrm{NaN}\ (\mathrm{NaN}\%)$ & & $\mathrm{NaN}\ (\mathrm{NaN}\%)$ \\
    & 49 & & $\mathrm{NaN}\ (\mathrm{NaN}\%)$ & & $\mathrm{NaN}\ (\mathrm{NaN}\%)$ & & $\mathrm{NaN}\ (\mathrm{NaN}\%)$ \\
    & 50 & & $\mathrm{NaN}\ (\mathrm{NaN}\%)$ & & $\mathrm{NaN}\ (\mathrm{NaN}\%)$ & & $\mathrm{NaN}\ (\mathrm{NaN}\%)$ \\
    & 51 & & $\mathrm{NaN}\ (\mathrm{NaN}\%)$ & & $\mathrm{NaN}\ (\mathrm{NaN}\%)$ & & $\mathrm{NaN}\ (\mathrm{NaN}\%)$ \\
    \midrule
    \multirow{3}{*}{Case V} & 10 & \multirow{3}{*}{75.2066147} & $\mathrm{NaN}\ (\mathrm{NaN}\%)$ & \multirow{3}{*}{6.0646468} & $\mathrm{NaN}\ (\mathrm{NaN}\%)$ & \multirow{3}{*}{0.025757} & $\mathrm{NaN}\ (\mathrm{NaN}\%)$ \\
    & 20 & & $\mathrm{NaN}\ (\mathrm{NaN}\%)$ & & $\mathrm{NaN}\ (\mathrm{NaN}\%)$ & & $\mathrm{NaN}\ (\mathrm{NaN}\%)$ \\
    & 40 & & $\mathrm{NaN}\ (\mathrm{NaN}\%)$ & & $\mathrm{NaN}\ (\mathrm{NaN}\%)$ & & $\mathrm{NaN}\ (\mathrm{NaN}\%)$ \\
    \midrule
    \multirow{5}{*}{Case VI} & 47 & \multirow{5}{*}{69.5002801} & $\mathrm{NaN}\ (\mathrm{NaN}\%)$ & \multirow{5}{*}{21.3088199} & $\mathrm{NaN}\ (\mathrm{NaN}\%)$ & \multirow{5}{*}{0.033759} & $\mathrm{NaN}\ (\mathrm{NaN}\%)$ \\
    & 48 & & $\mathrm{NaN}\ (\mathrm{NaN}\%)$ & & $\mathrm{NaN}\ (\mathrm{NaN}\%)$ & & $\mathrm{NaN}\ (\mathrm{NaN}\%)$ \\
    & 49 & & $\mathrm{NaN}\ (\mathrm{NaN}\%)$ & & $\mathrm{NaN}\ (\mathrm{NaN}\%)$ & & $\mathrm{NaN}\ (\mathrm{NaN}\%)$ \\
    & 50 & & $\mathrm{NaN}\ (\mathrm{NaN}\%)$ & & $\mathrm{NaN}\ (\mathrm{NaN}\%)$ & & $\mathrm{NaN}\ (\mathrm{NaN}\%)$ \\
    & 51 & & $\mathrm{NaN}\ (\mathrm{NaN}\%)$ & & $\mathrm{NaN}\ (\mathrm{NaN}\%)$ & & $\mathrm{NaN}\ (\mathrm{NaN}\%)$ \\
    \midrule
    \multirow{3}{*}{Case VII} & 95 & \multirow{3}{*}{70.4046528} & $\mathrm{NaN}\ (\mathrm{NaN}\%)$ & \multirow{3}{*}{6.2345658} & $\mathrm{NaN}\ (\mathrm{NaN}\%)$ & \multirow{3}{*}{0.032753} & $\mathrm{NaN}\ (\mathrm{NaN}\%)$ \\
    & 100 & & $\mathrm{NaN}\ (\mathrm{NaN}\%)$ & & $\mathrm{NaN}\ (\mathrm{NaN}\%)$ & & $\mathrm{NaN}\ (\mathrm{NaN}\%)$ \\
    & 105 & & $\mathrm{NaN}\ (\mathrm{NaN}\%)$ & & $\mathrm{NaN}\ (\mathrm{NaN}\%)$ & & $\mathrm{NaN}\ (\mathrm{NaN}\%)$ \\
    \bottomrule
  \end{tabular}%
  }
\end{table}
\begin{table}[htbp]
  \centering
  \caption{Almost Exact Simulation Results (Case I-VII)}
  \label{tab:terminal_part2}
  \resizebox{\textwidth}{!}{%
  \begin{tabular}{lclcccc}
    \toprule
    \multirow{2}{*}{Case} & \multirow{2}{*}{Strike} & \multicolumn{2}{c}{Almost Exact Simulation (Inverse Gaussian)} & \multicolumn{2}{c}{Almost Exact Simulation (Gamma)} \\
    \cmidrule(lr){3-4} \cmidrule(lr){5-6}
    & & Time(s) & Option Bias & Time(s) & Option Bias \\
    \midrule
    Case I & 1 & 2.6390925 & $0.001326\ (0.30\%)$ & 2.164342 & $0.0008747\ (0.20\%)$ \\
    Case I & 95 &  & $-8.215\ (-79.26\%)$ &  & $-7.956\ (-76.77\%)$ \\
    \midrule
    \multirow{3}{*}{Case II} & 100 & 0.2152987 & $-6.213\ (-84.12\%)$ & 0.204456 & $-5.994\ (-81.16\%)$ \\
    & 105 &  & $-4.361\ (-88.31\%)$ &  & $-4.195\ (-84.95\%)$ \\
    & 95 &  & $-11.64\ (-89.30\%)$ &  & $-11.97\ (-99.56\%)$ \\
    \midrule
    \multirow{3}{*}{Case III} & 100 & 0.256394 & $-8.96\ (-99.57\%)$ & 0.241701 & $-8.975\ (-99.73\%)$ \\
    & 105 &  & $-6.693\ (-99.74\%)$ &  & $-6.699\ (-99.84\%)$ \\
    & 95 &  & $-11.64\ (-99.30\%)$ &  & $-11.975\ (-99.56\%)$ \\
    \midrule
    \multirow{5}{*}{Case IV} & 47 & 20.5061735 & $-5.559\ (-78.96\%)$ & 13.86428 & $38.31\ (544.13\%)$ \\
    & 48 &  & $-5.287\ (-81.34\%)$ &  & $37.91\ (583.29\%)$ \\
    & 49 &  & $-5.135\ (-83.88\%)$ &  & $37.37\ (610.49\%)$ \\
    & 50 &  & $-4.903\ (-86.01\%)$ &  & $36.87\ (646.82\%)$ \\
    & 51 &  & $-4.664\ (-87.91\%)$ &  & $36.36\ (685.41\%)$ \\
    \midrule
    \multirow{3}{*}{Case V} & 10 & 22.739208 & $-5.255\ (-13.14\%)$ & 13.84294 & $43.83\ (109.58\%)$ \\
    & 20 &  & $-5.259\ (-17.53\%)$ &  & $43.83\ (146.08\%)$ \\
    & 40 &  & $-6.356\ (-53.29\%)$ &  & $41.99\ (352.08\%)$ \\
    \midrule
    \multirow{5}{*}{Case VI} & 47 & 0.1932632 & $-0.06061\ (-0.48\%)$ & 0.186914 & $-0.1563\ (-1.23\%)$ \\
    & 48 &  & $-0.05685\ (-0.46\%)$ &  & $-0.1489\ (-1.20\%)$ \\
    & 49 &  & $-0.05316\ (-0.44\%)$ &  & $-0.1419\ (-1.18\%)$ \\
    & 50 &  & $-0.04984\ (-0.42\%)$ &  & $-0.1354\ (-1.15\%)$ \\
    & 51 &  & $-0.04649\ (-0.41\%)$ &  & $-0.1291\ (-1.13\%)$ \\
    \midrule
    \multirow{3}{*}{Case VII} & 95 & 0.2169977 & $-11.64\ (-99.30\%)$ & 0.235242 & $-11.67\ (-99.56\%)$ \\
    & 100 &  & $-8.959\ (-99.57\%)$ &  & $-8.974\ (-99.73\%)$ \\
    & 105 &  & $-6.692\ (-99.74\%)$ &  & $-6.698\ (-99.84\%)$ \\
    \bottomrule
  \end{tabular}%
  }
\end{table}

As the table shows, the three ``exact simulation'' terminal methods were only able to output valid results in the single strike scenario of the Case I benchmark example. In all other test scenarios, they suffered complete numerical failure and do not possess full-scenario engineering practicality.

Specifically, the Baldeaux (2012) Exact Simulation method yielded valid pricing results only in the single strike scenario of Case I, where the relative bias approached $0$, confirming its theoretical unbiasedness. However, across the remaining 6 sets of test cases (Case II to Case VII) involving 22 strike price scenarios, it experienced numerical overflow and computational failure, outputting Not-a-Number (NaN) results. The core cause lies in the fact that the method requires the path-by-path calculation of high-order modified Bessel functions and inverse Laplace transforms. Under non-benchmark parameters, this easily triggers numerical overflow in the Bessel functions and integration convergence failures. It can only compute successfully under the literature's benchmark parameters, failing to adapt to the complex parameter environments of actual markets. Both the Fourier-Cosine Exact Simulation method and the Moment Matching Exact Simulation method behaved nearly identically to the Baldeaux method: they only succeeded in Case I's single strike scenario, and even there exhibited terrible pricing accuracy with relative biases of $-11.23\%$ and $-11.77\%$, respectively. They failed entirely in all other scenarios. The root causes of failure for these two methods are: for the Fourier-Cosine method, the cosine series truncation for the integrated variance distribution amplifies truncation error in the distribution tail under non-benchmark parameters, leading to failure in evaluating the characteristic function; for the moment matching method, matching only the first four moments via Gram-Charlier expansion cannot capture the non-normal heavy-tailed feature of the 3/2 model's variance distribution. This leads to non-physical negative densities, ultimately causing the inverse transform sampling of the cumulative distribution function to fail.

Regarding the two near-exact simulation methods, their pricing biases were only controllable under the medium-term, positive-correlation ATM scenario of Case VI. In the vast majority of other scenarios, they exhibited systematic and massive pricing errors. In the 5 near-ATM scenarios of Case VI, the relative bias of the IG method ranged from $-0.41\%$ to $-0.48\%$, and the Gamma method ranged from $-1.13\%$ to $-1.23\%$. This bias is manageable, and the computation time was roughly $0.19$ seconds, indicating high efficiency. However, in all other scenarios, both methods demonstrated significant pricing deviations: in the high vol-of-vol, strongly negative correlation scenarios (Cases II, III, VII), their relative biases reached $-80\%$ to $-99.8\%$ across all strike prices, completely deviating from the benchmark. In short-term, deep ITM scenarios (Cases IV, V), the Gamma method exhibited positive systematic biases exceeding $500\%$, and the IG method showed negative systematic biases of $-78\%$ to $-87\%$, losing all reference value. This happens because, under strong negative correlation, high vol-of-vol, and short-term scenarios, the distribution of the integrated variance becomes extremely heavy-tailed and highly right-skewed. The two-parameter IG and Gamma distributions, which only match the first two moments, cannot fit the higher-order features and tail behaviors of the distribution, resulting in systematic pricing errors.

\section{Result Analysis of Fourier Frequency-Domain Pricing Methods}

Frequency-domain pricing methods entirely bypass the path generation process of Monte Carlo simulations, solving for option prices directly through the Fourier integral of the characteristic function. This category encompasses the Fast Fourier Transform (FFT) and Fourier-Cosine (COS) methods. The complete results are detailed in Table \ref{tab:frequency}.

\begin{table}[htbp]
  \centering
  \caption{Fourier Frequency Domain Pricing Results}
  \label{tab:frequency}
  \resizebox{\textwidth}{!}{%
  \begin{tabular}{lclcccc}
    \toprule
    \multirow{2}{*}{Case} & \multirow{2}{*}{Strike} & \multicolumn{2}{c}{Fast Fourier Transformation} & \multicolumn{2}{c}{Fourier-Cosine Method} \\
    \cmidrule(lr){3-4} \cmidrule(lr){5-6}
    & & Time (s) & Option Bias & Time (s) & Option Bias \\
    \midrule
    Case I & 1 & 0.3843623 & $-2.434\times10^{-7}\ (-0.00\%)$ & 0.1115441 & $-2.437\times10^{-7}\ (-0.00\%)$ \\
    \midrule
    \multirow{3}{*}{Case II}
    & 95 & \multirow{3}{*}{0.1986538} & $-0.0004891\ (-0.00\%)$ & \multirow{3}{*}{0.0517513} & $-0.0004891\ (-0.00\%)$ \\
    & 100 & & $-0.0001461\ (-0.00\%)$ & & $-0.0001461\ (-0.00\%)$ \\
    & 105 & & $-0.0009548\ (-0.02\%)$ & & $-0.0009548\ (-0.02\%)$ \\
    \midrule
    \multirow{3}{*}{Case III}
    & 95 & \multirow{3}{*}{0.2813563} & $-0.0005218\ (-0.00\%)$ & \multirow{3}{*}{0.0735898} & $-0.0005218\ (-0.00\%)$ \\
    & 100 & & $-0.001246\ (-0.01\%)$ & & $-0.001246\ (-0.01\%)$ \\
    & 105 & & $-0.0008511\ (-0.01\%)$ & & $-0.0008511\ (-0.01\%)$ \\
    \midrule
    \multirow{5}{*}{Case IV}
    & 47 & \multirow{5}{*}{15.8604601} & $0.0003048\ (0.00\%)$ & \multirow{5}{*}{3.7037431} & $0.0003048\ (0.00\%)$ \\
    & 48 & & $0.068\ (1.05\%)$ & & $0.068\ (1.05\%)$ \\
    & 49 & & $0.001088\ (0.02\%)$ & & $0.001088\ (0.02\%)$ \\
    & 50 & & $0.001157\ (0.02\%)$ & & $0.001157\ (0.02\%)$ \\
    & 51 & & $0.001138\ (0.02\%)$ & & $0.001138\ (0.02\%)$ \\
    \midrule
    \multirow{3}{*}{Case V}
    & 10 & \multirow{3}{*}{16.0524815} & $0.0015\ (0.00\%)$ & \multirow{3}{*}{3.788367} & $0.0015\ (0.00\%)$ \\
    & 20 & & $0.001557\ (0.01\%)$ & & $0.001557\ (0.01\%)$ \\
    & 40 & & $0.002021\ (0.02\%)$ & & $0.002021\ (0.02\%)$ \\
    \midrule
    \multirow{5}{*}{Case VI}
    & 47 & \multirow{5}{*}{0.2070812} & $0.01801\ (0.14\%)$ & \multirow{5}{*}{0.0686066} & $0.01785\ (0.14\%)$ \\
    & 48 & & $0.01801\ (0.15\%)$ & & $0.01791\ (0.15\%)$ \\
    & 49 & & $0.01807\ (0.15\%)$ & & $0.01791\ (0.15\%)$ \\
    & 50 & & $0.01794\ (0.15\%)$ & & $0.01778\ (0.15\%)$ \\
    & 51 & & $0.01794\ (0.16\%)$ & & $0.01775\ (0.16\%)$ \\
    \midrule
    \multirow{3}{*}{Case VII}
    & 95 & \multirow{3}{*}{0.3207817} & $-2.18\times10^{-5}\ (-0.00\%)$ & \multirow{3}{*}{0.0993632} & $-2.179\times10^{-5}\ (-0.00\%)$ \\
    & 100 & & $-4.598\times10^{-5}\ (-0.00\%)$ & & $-4.597\times10^{-5}\ (-0.00\%)$ \\
    & 105 & & $4.895\times10^{-5}\ (0.00\%)$ & & $4.895\times10^{-5}\ (0.00\%)$ \\
    \bottomrule
  \end{tabular}%
  }
\end{table}

As the table shows, these two methods are the only ones in this experiment that experienced zero numerical failures across all 7 test cases and all strike prices, while simultaneously maintaining exceptionally high pricing accuracy. Their comprehensive performance is significantly superior to all Monte Carlo-based methods.

In terms of numerical stability, the FFT and COS methods never produced non-numerical or infinite results in any scenario. Whether evaluating benchmark ATM scenarios, extreme high vol-of-vol and strong negative correlation cases, short maturity scenarios, or deep ITM/OTM scenarios, they stably output pricing results, entirely circumventing the sampling failures and distribution fitting failures associated with Monte Carlo methods. In terms of pricing accuracy, the relative biases of both methods were contained within $1.05\%$ across all scenarios, with the vast majority falling below $0.02\%$. In the extreme high vol-of-vol environment of Case II, their maximum relative bias was merely $-0.02\%$. In the short maturity Case IV, the maximum bias was $1.05\%$, while for all other strike prices it remained under $0.02\%$. In the high vol-of-vol Case VI, the relative bias stood around $0.15\%$. This precision overwhelmingly surpasses all Monte Carlo methods. The core reason is that frequency-domain methods execute pricing directly via numerical integration relying on the closed-form solution of the 3/2 model's characteristic function. They are immune to Monte Carlo sampling errors, distribution fitting errors, and time-discretization truncation errors, suffering only from negligible numerical integration truncation errors.

Regarding computational efficiency, the COS method comprehensively outperformed the FFT method. Across all test scenarios, the execution time for the COS method was only $1/3$ to $1/4$ of that required by FFT. In the multi-strike Case IV, FFT took 15.86 seconds, whereas COS took a mere 3.70 seconds; in single-strike scenarios, COS execution times ranged from 0.05 to 0.11 seconds, boasting efficiency far superior to any Monte Carlo scheme. The theoretical driver behind this efficiency gap is that the COS method, based on cosine series expansion, converges exponentially. It requires only 128 truncation terms to match the precision of an FFT grid with 4096 points. Moreover, the COS method is free from the FFT's picket-fence effect, allowing flexible pricing at arbitrary strike prices without extra interpolation steps, thereby further reducing computational overhead.

\section{Comprehensive Performance Comparison and Research Conclusions}

Based on the full-scenario experimental results from the 7 test cases, the 9 classes of pricing methods can be systematically tiered across three core dimensions: numerical stability, pricing accuracy, and computational efficiency. 

The first tier consists of the COS and FFT frequency-domain methods, which are the only approaches that ensure full-scenario numerical stability while retaining extreme accuracy and efficiency. Specifically, the COS method comprehensively outperforms the FFT method, cementing its status as the optimal engineering solution for pricing European options under the 3/2 model. 

The second tier includes the Exact Stepping, Poisson-Gamma Near-Exact Stepping, and QE Near-Exact Stepping methods from the time-stepping category. They represent the only stable solutions across all scenarios within the Monte Carlo framework. They can stably generate complete variance paths and offer a favorable balance between pricing accuracy and efficiency, making them the sole reliable options for pricing path-dependent options or for dynamic hedging scenarios where complete path generation is mandatory. 

The third tier is the Euler/Milstein discretization scheme, which is only valid under fine step sizes and conventional parameters. It fails numerically in extreme scenarios and exhibits extremely poor efficiency; thus, it serves merely as a baseline for theoretical accuracy comparisons. 

The fourth tier comprises the IG and Gamma Near-Exact Terminal Simulation methods, which provide controllable bias only in specific ATM scenarios. In the overwhelming majority of cases, they suffer massive systematic pricing errors. They are limited to specific academic comparisons and possess no engineering practicality. 

The fifth tier features the Baldeaux Exact Simulation, Fourier-Cosine Exact Terminal Simulation, and Moment Matching Exact Terminal Simulation methods. These function solely in a singular benchmark scenario, utterly failing in all others. They are strictly limited to validating theoretical unbiasedness.

Based on these full-scenario experimental validations, three conclusions can be drawn. First, the theoretical unbiasedness of a 3/2 model pricing method does not equate to full-scenario stability in engineering applications. Nominally ``exact'' terminal simulation methods are highly susceptible to special function overflows and integration convergence failures under non-benchmark parameters. They only pass theoretical validation under literature-provided benchmarks and fail to adapt to the complex parameters of actual markets. Second, time-stepping sampling methods built upon the distributional properties of the variance process constitute the optimal implementation schemes within the Monte Carlo framework. They inherently guarantee the non-negativity of the variance process, ensuring numerical stability in all extreme environments, and consistently outperform traditional Euler discretization in accuracy, providing a robust path-generation basis for derivatives pricing. Third, the characteristic function-based COS frequency-domain method is the absolute optimal choice for pricing European options under the 3/2 model. It preserves supreme precision and stability across all test scenarios with zero numerical failures, and its computational efficiency far exceeds all Monte Carlo techniques, granting it overwhelming performance superiority in batch option pricing, model parameter calibration, and extreme market conditions.

Deriving from this holistic validation, we provide corresponding method recommendations for distinct application scenarios. For batch pricing of European options and model parameter calibration, the COS frequency-domain method is the premier choice; its stability, peak precision, and supreme efficiency seamlessly satisfy these demands. For scenarios necessitating complete variance paths, such as path-dependent option pricing or dynamic hedging, the Poisson-Gamma or QE Near-Exact Stepping methods are prioritized, given their comprehensive stability and well-balanced precision and efficiency. For baseline theoretical precision comparisons, the Baldeaux Exact Simulation method may be utilized, but strictly under benchmark parameter settings. The remaining terminal simulation methods are not recommended for practical pricing environments and should be reserved solely for specific academic comparative analyses.